% TEMPLATE-MILC-v3.tex - plantilla maestra para todas las guias MILC v3
% Ver scripts/generadores/build-guia-milc.py para la lista completa de
% claves esperadas. El script valida que no queden placeholders sin
% reemplazar antes de escribir el archivo final.

\documentclass[11pt,letterpaper]{article}
\usepackage[margin=1.75cm]{geometry}
\usepackage[table]{xcolor}
\usepackage{fontspec}
\usepackage[spanish]{babel}
\usepackage{tikz}
\usepackage[most]{tcolorbox}
\usepackage{tabularx}
\usepackage{array}
\usepackage{fancyhdr}
\usepackage{titlesec}
\usepackage{ragged2e}
\usepackage{enumitem}
\usepackage{microtype}

\setmainfont{Helvetica Neue}
\setsansfont{Helvetica Neue}
\setlength{\parindent}{0pt}
\setlength{\parskip}{6pt}
\hyphenpenalty=200
\exhyphenpenalty=200
\pretolerance=400
\tolerance=2000
\setlength{\emergencystretch}{1em}
\renewcommand{\arraystretch}{1.25}
\setlist{leftmargin=5mm,itemsep=1mm,topsep=1mm}
\newcolumntype{Y}{>{\RaggedRight\arraybackslash}X}

\definecolor{milcMagenta}{HTML}{E6007E}
\definecolor{milcVino}{HTML}{5A0038}
\definecolor{milcTurquesa}{HTML}{0093A5}
\definecolor{milcVerde}{HTML}{58A923}
\definecolor{milcMostaza}{HTML}{E5B400}
\definecolor{milcOcre}{HTML}{B86600}
\definecolor{milcNegro}{HTML}{241816}
\definecolor{milcGris}{HTML}{F7F7F4}
\definecolor{milcLinea}{HTML}{E8E2DD}
\definecolor{milcGradePrimary}{HTML}{0066FF}
\definecolor{milcGradeDark}{HTML}{003D99}
\definecolor{milcGradeSoft}{HTML}{D6E8FF}
\definecolor{milcGradeAccent}{HTML}{CFE7FF}
\definecolor{milcGradeText}{HTML}{FFFFFF}
\color{milcNegro}

\pagestyle{fancy}
\fancyhf{}
\lhead{\small Tecnología e Informática · Grado 11}
\rhead{\small Guía 6 · MILC v3}
\cfoot{\small\thepage}
\renewcommand{\headrulewidth}{0pt}

\newtcolorbox{titlebox}[1]{
  enhanced,colback=#1,colframe=#1,arc=7mm,boxrule=0pt,
  left=7mm,right=7mm,top=3mm,bottom=3mm,
  before skip=8pt,after skip=8pt,
  fontupper=\sffamily\bfseries\Large\color{white}
}

\newtcolorbox{softbox}[3]{
  enhanced,colback=#2,colframe=#1,borderline west={4pt}{0pt}{#1},
  arc=2mm,boxrule=.4pt,left=4mm,right=4mm,top=3mm,bottom=3mm,
  before skip=6pt,after skip=8pt,title={#3},coltitle=white,
  colbacktitle=#1,fonttitle=\sffamily\bfseries,fontupper=\RaggedRight,
  attach boxed title to top left={xshift=3mm,yshift=-2mm},
  boxed title style={arc=2mm, boxrule=0pt}
}

\newtcolorbox{drawbox}[2]{
  enhanced,colback=white,colframe=#1,arc=4mm,boxrule=1pt,
  left=5mm,right=5mm,top=4mm,bottom=4mm,before skip=8pt,after skip=8pt,
  title={#2},coltitle=white,colbacktitle=#1,fonttitle=\sffamily\bfseries,
  fontupper=\RaggedRight,
  attach boxed title to top left={xshift=4mm,yshift=-2mm},
  boxed title style={arc=3mm, boxrule=0pt}
}

\newtcolorbox{infoband}[1]{
  enhanced,colback=#1!8,colframe=#1!50,arc=2mm,boxrule=.5pt,
  left=4mm,right=4mm,top=2mm,bottom=2mm,before skip=4pt,after skip=4pt,
  fontupper=\small\RaggedRight
}

% Puente narrativo entre secciones (contrato editorial Conéctate).
% Da continuidad: "Acabas de X. Ahora vas a Y porque Z."
% Visualmente: banda izquierda en color del grado, texto en itálica pequeña.
\newtcolorbox{puentebox}{
  enhanced,colback=milcGradeSoft,colframe=milcGradeDark,
  borderline west={3pt}{0pt}{milcGradeDark},
  arc=0mm,boxrule=0pt,left=5mm,right=4mm,top=2.5mm,bottom=2.5mm,
  before skip=4pt,after skip=8pt,
  fontupper=\itshape\small\color{milcNegro}\RaggedRight
}
\newcommand{\puente}[1]{%
  \begin{puentebox}%
    \textcolor{milcGradeDark}{\textbf{\textrightarrow}}\hspace{2mm}#1%
  \end{puentebox}%
}

\newcommand{\linead}[1]{\noindent\color{milcLinea}\rule{#1}{0.5pt}\color{milcNegro}}
\newcommand{\respuesta}[1]{\par\vspace{2mm}\foreach \n in {1,...,#1}{\noindent\color{milcLinea}\rule{\linewidth}{0.5pt}\par\vspace{5mm}}\color{milcNegro}}
\newcommand{\checkbox}{\textcolor{milcGradeDark}{\fbox{\phantom{x}}}\hspace{5pt}}

\newcommand{\phasecard}[4]{
\begin{tcolorbox}[enhanced,colback=#3,colframe=#2,arc=3mm,boxrule=.5pt,height=3.05cm,valign=center,halign=center,left=1.5mm,right=1.5mm,top=2mm,bottom=2mm]
{\sffamily\bfseries\fontsize{22}{24}\selectfont\textcolor{#2}{#1}}\par
{\sffamily\bfseries\small #4}
\end{tcolorbox}
}

\begin{document}
\thispagestyle{empty}
\begin{tikzpicture}[remember picture,overlay]
  \fill[white] (current page.south west) rectangle (current page.north east);
  \path[fill=milcGradePrimary,rounded corners=16mm]
    ([xshift=.55cm,yshift=-.55cm]current page.north west) rectangle
    ([xshift=-.55cm,yshift=3.65cm]current page.south east);

  \node[anchor=north west,text=milcGradeText] at ([xshift=2.0cm,yshift=-1.85cm]current page.north west)
    {\sffamily\bfseries\fontsize{14}{16}\selectfont GRADO};
  \node[anchor=north west,text=milcGradeText,opacity=.82] at ([xshift=2.0cm,yshift=-2.75cm]current page.north west)
    {\sffamily\bfseries\fontsize{9}{11}\selectfont SERIE GUÍAS · TECNOLOGÍA E INFORMÁTICA};

  \begin{scope}[shift={([xshift=-3.05cm,yshift=-2.55cm]current page.north east)}]
    \fill[white,opacity=.80] (-.62,-.52) rectangle (.62,.52);
    \draw[milcGradeText,line width=1pt] (-.62,-.52) rectangle (.62,.52);
    \fill[milcGradeDark] (-.42,-.38) rectangle (-.22,.06);
    \fill[milcGradeAccent] (-.10,-.38) rectangle (.10,.24);
    \fill[milcGradePrimary!60!black] (.22,-.38) rectangle (.42,.38);
  \end{scope}

  \node[anchor=west,text=milcGradeText] at ([xshift=1.9cm,yshift=-12.85cm]current page.north west)
    {\sffamily\bfseries\fontsize{124}{124}\selectfont 11};
  \draw[milcGradeText,line width=1.7pt]
    ([xshift=4.52cm,yshift=-11.68cm]current page.north west) circle (.13cm);

  \node[anchor=west,text=milcGradeText,text width=8.7cm,align=left] at ([xshift=10.35cm,yshift=-11.6cm]current page.north west)
    {
     {\sffamily\bfseries\fontsize{19}{22}\selectfont SEO básico ---\\hacerme visible\\sin engaños}\\[4mm]
     {\sffamily\bfseries\fontsize{23}{26}\selectfont Guía 6-11-TIC}\\[2mm]
     {\sffamily\fontsize{13}{16}\selectfont Período 1 · Presencia y marca digital}\\[5mm]
     {\sffamily\bfseries\fontsize{11}{13}\selectfont Institución Educativa Sor María Juliana}\\[1mm]
     {\sffamily\fontsize{10}{12}\selectfont Autor: Dr. Álvaro Cárdenas Orozco}
    };

  \path[fill=milcGradeSoft,rounded corners=7mm]
    ([xshift=1.35cm,yshift=.85cm]current page.south west) rectangle
    ([xshift=-1.35cm,yshift=4.25cm]current page.south east);
  \draw[milcGradeDark,line width=.65pt,rounded corners=7mm]
    ([xshift=1.35cm,yshift=.85cm]current page.south west) rectangle
    ([xshift=-1.35cm,yshift=4.25cm]current page.south east);
  \node[anchor=center,text=milcNegro,text width=17.0cm,align=center] at ([yshift=2.55cm]current page.south)
    {\sffamily\fontsize{10}{12}\selectfont Metodología MILC · Apertura ancestral · Pensamiento computacional · Triángulo Dussel-Estoicismo-Floridi\\
    \textcolor{milcGradeDark}{\bfseries Producto:} Plan de palabras clave + meta-etiquetas configuradas + sitemap.xml del sitio};
\end{tikzpicture}
\mbox{}
\newpage

\section*{Apertura: SEO básico --- hacerme visible sin engaños}

\begin{softbox}{milcGradeDark}{milcGradeSoft}{Saber ancestral}
En las plazas del Valle, el \textbf{pregonero} era el oficio de hacer llegar el mensaje correcto a la gente correcta. No gritaba: \emph{pregonaba}. Anunciaba con voz clara, a la hora justa, en el lugar donde lo necesitaban escuchar. ``Los huevos del corral de Don Tobías'', ``el agua bendita del padre Antonio'', ``el pan tibio de la señora Inés''. El que lo oía iba; el que no lo necesitaba seguía su camino. Pregón honesto, no grito desesperado.
\end{softbox}

\begin{softbox}{milcTurquesa}{milcGris}{Saber contemporáneo}
El \textbf{SEO} (Search Engine Optimization) es el pregón digital. No truquillos para engañar a Google: técnica honesta para que tu sitio aparezca cuando alguien busca lo que tú ofreces. Tres piezas: \textbf{palabras clave} (qué busca la gente), \textbf{meta-etiquetas} (cómo te presentas en el resultado), y \textbf{sitemap} (el mapa que le entregas al buscador).
\end{softbox}

\begin{softbox}{milcMostaza}{milcGris}{Pregunta-puente}
\textit{¿Qué hacía que el pregonero antiguo lograra que la gente \emph{correcta} lo escuchara entre el ruido de la plaza? ¿Por qué a algunos los oían y a otros no?}
\end{softbox}

\begin{softbox}{milcVerde}{milcGris}{Saber y saber hacer}
Al terminar podrás: (1) \textbf{identificar} las palabras clave que la gente buscaría para encontrar a alguien como tú; (2) \textbf{explicar} con tus palabras cómo decide un buscador qué mostrar primero; (3) \textbf{aplicar} meta-etiquetas, sitemap y alt text en tu sitio; (4) \textbf{evaluar} si apareces al buscarte.
\end{softbox}

\small
\begin{tabularx}{\textwidth}{>{\bfseries}p{3.0cm}Y}
\rowcolor{milcGradePrimary!12}Autor & Dr. Álvaro Cárdenas Orozco\\
DBA & Comunicación profesional en entornos digitales (MEN, Lineamientos T\&I)\\
Referentes & Floridi (infoética) · Dussel (estética de la liberación) · Estoicismo\\
Producto final & Plan de palabras clave + meta-etiquetas configuradas + sitemap.xml del sitio\\
\end{tabularx}
\normalsize

\puente{Antes de empezar, mira el viaje completo. Hoy le abrirás la \emph{voz} a tu sitio --- haces que Google y otros buscadores lo encuentren cuando alguien busca lo que tu marca ofrece.}

\begin{titlebox}{milcGradeDark}
Ruta MILC: así vamos a aprender
\end{titlebox}
\begin{tabularx}{\textwidth}{>{\centering\arraybackslash}X>{\centering\arraybackslash}X>{\centering\arraybackslash}X>{\centering\arraybackslash}X}
\phasecard{1}{milcVerde}{milcGris}{Escuta\\[-1mm]\small Observo, escucho y pregunto.} &
\phasecard{2}{milcTurquesa}{milcGris}{Sistematización\\[-1mm]\small Pensamiento computacional.} &
\phasecard{3}{milcMagenta}{milcGris}{Praxis\\[-1mm]\small Construyo y aplico.} &
\phasecard{4}{milcVino}{milcGris}{Evaluación\\[-1mm]\small Reflexión liberadora.}
\end{tabularx}


\puente{Antes de optimizar tu sitio para Google, hay que entender cómo busca la gente real. Vas a ponerte en los zapatos de quien podría buscarte: ¿qué palabras escribiría en Google si nunca te ha oído nombrar?}

\begin{titlebox}{milcVerde}
Fase 1 · Escuta
\end{titlebox}

\begin{softbox}{milcVerde}{milcGris}{Escena para escuchar}
\textbf{Actividad 1 · IDENTIFICA --- Palabras clave del que te busca} (15 min · individual). Imagina que alguien (futuro empleador, cliente, profesor) quiere encontrar a alguien como tú pero no te conoce. Anota \textbf{10 frases o palabras} que esa persona escribiría en Google. No tus términos técnicos: los términos que usaría alguien que apenas sabe lo que necesita.
\end{softbox}

\checkbox Pensé en al menos 3 perfiles distintos de personas que podrían buscarme.\\
\checkbox Anoté frases largas (3+ palabras) no solo términos sueltos.\\
\checkbox Probé en Google una de esas frases para ver qué aparece hoy.

\begin{infoband}{milcVerde}
\textbf{Lo que va al cuaderno (Actividad 1 · IDENTIFICA).} Título: «Actividad 1 --- Palabras clave de mi marca». \textbf{Formato:} lista numerada de 10 frases de búsqueda, agrupadas en 3 perfiles imaginados (quién busca cada frase). Después un párrafo de 3 renglones eligiendo las 3 frases más estratégicas. \textbf{Sabes que terminaste cuando} las frases suenan a búsquedas reales (3+ palabras, lenguaje cotidiano, no jerga técnica) y las 3 elegidas son las que tu sitio podría sostener honestamente.
\end{infoband}

\puente{Ya pensaste como quien busca. Ahora vamos a entender \textbf{cómo decide Google} qué mostrar primero. No es magia ni suerte: hay reglas claras y un puñado de decisiones técnicas que sí dependen de ti.}

\begin{titlebox}{milcTurquesa}
Fase 2 · Sistematización con pensamiento computacional
\end{titlebox}

\begin{softbox}{milcTurquesa}{milcGris}{Cuatro pilares aplicados al tema}
Google decide qué mostrar primero según un puñado de señales que tú puedes controlar. Vamos a descomponerlas con los pilares del pensamiento computacional.
\end{softbox}

\small
\begin{tabularx}{\textwidth}{>{\bfseries\RaggedRight\arraybackslash}p{3.6cm}Y}
\rowcolor{milcTurquesa!90}\color{white}Pilar & \color{white}\bfseries Aplicación al tema\\
1. Descomposición & \textbf{Descomposición.} El SEO se descompone en 4 capas que cada buscador evalúa: (1) \textbf{contenido} (las palabras reales de tu sitio --- responde a lo que la gente busca con texto sustantivo, no relleno); (2) \textbf{metadatos} (title de 50-60 caracteres + description de 140-160 caracteres que Google muestra como resultado --- son tu cartel en la plaza); (3) \textbf{estructura} (HTML semántico, jerarquía clara de encabezados h1 → h2 → h3, navegación predecible); (4) \textbf{accesibilidad} (alt text en imágenes, contraste WCAG, sitio móvil-friendly). Las 4 capas pesan; ignorar una baja el ranking.\\
2. Reconocimiento de patrones & \textbf{Reconocimiento de patrones.} Los buscadores siguen un patrón consistente desde Google PageRank hasta hoy: dan más peso a \textbf{coincidencia + autoridad + experiencia (E-E-A-T)}. \textbf{Coincidencia}: tu sitio menciona lo que se buscó en lenguaje cercano al del que busca. \textbf{Autoridad}: otros sitios reputados te enlazan (link building honesto). \textbf{Experiencia}: la gente se queda en tu sitio, navega, no rebota a los 3 segundos. Los algoritmos de Google premian estos 3 factores combinados; un sitio rica en uno solo no escala.\\
3. Abstracción & \textbf{Abstracción.} Lo esencial del SEO cabe en una pregunta: \emph{¿le facilito a Google entender de qué se trata mi sitio en 5 segundos?}. Si la respuesta es sí, vas bien. Si el bot tiene que adivinar (title genérico, h1 ambiguo, descripción vacía, imágenes sin alt), tu posición caerá sin importar cuánto contenido tengas. El SEO bueno hace lo opuesto al SEO de los noventa (que escondía palabras clave); hoy, claridad gana sobre truco.\\
4. Algoritmo conceptual & \textbf{Algoritmo.} (1) Elige 3 frases-clave estratégicas que combinen tu nombre + tu propuesta de valor; (2) ponlas en tu title y description sin sonar a relleno; (3) úsalas naturalmente en H1, H2 y párrafos --- una vez es suficiente, repetirlas 20 veces es \emph{keyword stuffing} y te penaliza; (4) escribe alt text en cada imagen describiendo lo que muestra (sirve a lectores de pantalla y a Google); (5) genera \texttt{sitemap.xml} y subelo a la raíz; (6) verifica todo en search.google.com/search-console.\\
\end{tabularx}
\normalsize

\begin{drawbox}{milcTurquesa}{Anatomía de una página optimizada para buscadores}
\textbf{Title (50-60 caracteres):} ``Tu nombre --- tu propuesta de valor''. Aparece en la pestaña y en el resultado de Google. \\[1mm] \textbf{Meta description (140-160 caracteres):} resumen claro de qué encontrará quien entre. \\[1mm] \textbf{H1 único:} con tu palabra clave principal de forma natural. \\[1mm] \textbf{Alt text en imágenes:} describe lo que la imagen muestra (sirve a lectores de pantalla \emph{y} a Google). \\[1mm] \textbf{Sitemap.xml:} lista de todas tus páginas, en la raíz de tu dominio.
\end{drawbox}

\begin{softbox}{milcMostaza}{milcGris}{Errores comunes y su corrección}
(1) Title genérico (``Inicio'' o ``Home''): pierdes el 90\% del SEO porque el title es la señal más fuerte que Google lee primero. (2) Repetir la palabra clave 20 veces en el contenido (\emph{keyword stuffing}): Google lo detecta desde 2012 y lo penaliza. (3) Imágenes sin alt text: invisibles para Google y para lectores de pantalla --- doble penalización ética y de SEO. (4) Promesas en el title o description que tu sitio no cumple: la gente entra, sale rápido (bounce rate alto) y tu posición cae en cuestión de semanas. (5) No tener versión móvil bien hecha: Google indexa móvil primero desde 2018; si tu sitio se rompe a 360 px, es como si no existiera.
\end{softbox}

\begin{infoband}{milcTurquesa}
\textbf{Lo que va al cuaderno (Actividad 2 · EXPLICA).} Título: «Actividad 2 --- Cómo decide Google qué mostrar primero». \textbf{Formato:} 4 fichas, una por capa del SEO (contenido / metadatos / estructura / accesibilidad), cada ficha con: qué es (1 frase con tus palabras) + un ejemplo concreto de cómo lo aplicarías a tu sitio. \textbf{Sabes que terminaste cuando} puedes explicar cada capa sin mirar notas y dar un ejemplo de tu propio caso.
\end{infoband}

\puente{Conoces el principio. Toca aplicarlo a \emph{tu} sitio: meta-etiquetas, sitemap, alt text en imágenes. Cinco cambios pequeños que multiplican tu visibilidad sin ningún engaño.}

\begin{titlebox}{milcMagenta}
Fase 3 · Praxis con pensamiento computacional
\end{titlebox}

\begin{softbox}{milcMagenta}{milcGris}{Construyendo el producto con los cuatro pilares}
\textbf{Actividad 3 · APLICA --- Optimizar tu sitio en 5 pasos} (35 min · individual). Tienes el principio. Hora de ejecutar en \emph{tu} sitio: meta-tags, alt text en imágenes, jerarquía de encabezados y sitemap. Al final, tu sitio estará listo para que Google lo encuentre cuando alguien busque lo que ofreces honestamente.
\end{softbox}

\small
\begin{tabularx}{\textwidth}{>{\bfseries\RaggedRight\arraybackslash}p{3.6cm}Y}
\rowcolor{milcMagenta!90}\color{white}Pilar & \color{white}\bfseries Aplicación al producto\\
Descomposición & \textbf{Descomposición.} Tu optimización se descompone en 5 cambios atómicos y secuenciales: (1) \textbf{title} de 50-60 caracteres específico; (2) \textbf{meta description} de 140-160 caracteres clara; (3) \textbf{H1 único} con tu palabra clave principal de forma natural; (4) \textbf{alt text} descriptivo en todas las imágenes; (5) \textbf{sitemap.xml} en la raíz del dominio. Hechos en orden y verificados uno a uno, son la base del SEO técnico básico para cualquier sitio personal.\\
Patrones & \textbf{Patrones.} Sigue el patrón \textbf{``honestidad SEO''}: cada palabra clave que pongas debe estar sostenida por contenido real y verificable. Si prometes en el title ``programador full stack'' pero tu sitio no tiene proyectos de programación visibles, Google te baja en cuestión de semanas porque el comportamiento de los visitantes (rebote rápido) le dice al algoritmo que no cumples. La regla de oro: SEO sin contenido honesto detrás es como un cartel sin tienda.\\
Abstracción & \textbf{Abstracción.} Cada meta-etiqueta expresa una sola idea esencial. Si tu description dice 3 cosas distintas (``soy programador, también diseño, también escribo''), Google la corta y muestra un fragmento aleatorio --- pierdes la primera impresión. Mejor una description que dice una cosa con claridad: ``Diseño sitios web para pequeños negocios de Cartago. 12 proyectos, 3 años de oficio''. Una idea, datos concretos, ninguna inflación.\\
Algoritmo de armado & \textbf{Algoritmo.} (1) Abre tu sitio en VS Code; (2) edita \texttt{<title>} y \texttt{<meta name="description">} en el \texttt{<head>} de \texttt{index.html}; (3) revisa que existe UN \texttt{<h1>} con tu palabra clave usada con naturalidad (no forzada); (4) recorre todas las imágenes y agrega \texttt{alt="descripción específica"}; (5) crea \texttt{sitemap.xml} con la lista de tus URLs; (6) commit y push; (7) sube el sitemap a Google Search Console para que lo indexe rápido.\\
\end{tabularx}
\normalsize

\begin{drawbox}{milcMagenta}{Checklist antes de declarar el SEO básico hecho}
\checkbox Title de 50-60 caracteres con tu propuesta de valor. \\[1mm] \checkbox Description de 140-160 caracteres clara y específica. \\[1mm] \checkbox H1 único con tu palabra clave principal de forma natural. \\[1mm] \checkbox Todas las imágenes tienen \texttt{alt} descriptivo. \\[1mm] \checkbox \texttt{sitemap.xml} en la raíz, validado en search.google.com/test.
\end{drawbox}

\begin{softbox}{milcMagenta}{milcGris}{Plantilla de trabajo}
\textbf{Title.} \textit{Tu fórmula final, 50-60 caracteres.} \\[1mm] \textbf{Description.} \textit{Tu resumen final, 140-160 caracteres.} \\[1mm] \textbf{Palabras clave estratégicas.} \textit{Las 3 elegidas, dónde aparecen en el sitio.} \\[1mm] \textbf{Sitemap.} \textit{URL pública del sitemap.xml.}
\end{softbox}

\begin{infoband}{milcMagenta}
\textbf{Lo que va al cuaderno (Actividad 3 · APLICA).} Título: «Actividad 3 --- Mi SEO básico aplicado». \textbf{Formato:} 4 secciones (Title · Description · Palabras clave aplicadas · Sitemap), cada una con sus datos finales y captura del código real. \textbf{Extensión:} una página de cuaderno. \textbf{Sabes que terminaste cuando} al buscar tu title en Google (búsqueda con comillas) tu sitio aparece en los primeros resultados.
\end{infoband}

\puente{Lo que sigue es el resumen de \emph{qué} entregas hoy. El criterio simple: que al buscar tu nombre o tu propuesta de valor en Google, tu sitio aparezca en la primera página.}

\begin{titlebox}{milcMostaza}
Producto de la sesión
\end{titlebox}
\textbf{Sitio optimizado con title, description, alt text y sitemap.xml}.
\begin{softbox}{milcMostaza}{milcGris}{Criterios de calidad}
(1) Title específico (no ``Inicio''), 50-60 caracteres. (2) Description única, 140-160 caracteres, sin promesas que no cumples. (3) UN H1 con palabra clave principal usada naturalmente. (4) Cien por ciento de imágenes con alt descriptivo. (5) Sitemap.xml accesible y validado.
\end{softbox}

\puente{Antes de cerrar, mira tu trabajo desde las cinco dimensiones humanas. El SEO es técnico pero tiene cara ética: invitar o engañar, servir o saturar.}

\begin{titlebox}{milcVino}
Fase 4 · Evaluación liberadora — cinco dimensiones
\end{titlebox}

\begin{softbox}{milcVino}{milcGris}{Sentido de la evaluación}
Evaluar no es solo revisar una nota. Es reconocer qué aprendí, qué cambió en mí, qué emociones manejé, cómo me relaciono con la ciudad y la comunidad, y qué le dejaré a quien viene detrás.
\end{softbox}

\textbf{Mi reflexión liberadora en cinco dimensiones.} Completa cada frase iniciada:

\small
\begin{tabularx}{\textwidth}{>{\bfseries\RaggedRight\arraybackslash}p{3.4cm}Y>{\RaggedRight\arraybackslash}p{4.6cm}}
\rowcolor{milcVino}\color{white}Dimensión & \color{white}\bfseries Mi respuesta (frame starter) & \color{white}\bfseries Algo que descubrí\\
1. Personal & Lo que más cambió en mí fue \linead{6.5cm} & \linead{4.2cm}\\
2. Emocional & La emoción más fuerte fue \linead{2.5cm} y la regulé \linead{3cm} & \linead{4.2cm}\\
3. Ciudadana & En democracia esto importa porque \linead{6cm} & \linead{4.2cm}\\
4. Local (Cartago) & En mi barrio sería útil para \linead{6.5cm} & \linead{4.2cm}\\
5. Intergeneracional & A mi abuela/abuelo le diría \linead{6.5cm} & \linead{4.2cm}\\
\end{tabularx}
\normalsize

\begin{infoband}{milcVino}
\textbf{Para la próxima sustentación.} Vuelve a esta tabla en seis meses. Verás que las respuestas cambian — no porque lo aprendido cambie, sino porque tú cambias. La evaluación liberadora es una conversación contigo a través del tiempo.
\end{infoband}


\puente{Y para terminar, tres voces sobre la visibilidad. Dussel sobre incluir vs promocionarse, Marco Aurelio sobre los algoritmos que cambian, Floridi sobre la búsqueda como acto cooperativo.}

\begin{titlebox}{milcGradeDark}
Triángulo de pensamiento · cierre filosófico
\end{titlebox}

\begin{softbox}{milcVino}{milcGris}{Dussel · lente del nosotros}
\textit{``Toda comunidad de comunicación se hace plena cuando incluye la voz del excluido.''}

El SEO es la diferencia entre que tu sitio sea encontrado por la gente que de verdad lo necesita o que se pierda en la marea de internet. Pero también puede ser usado para empujar voces hacia abajo: cuando un negocio grande compra términos que podrían servirle a uno pequeño del barrio. La pregunta ética está en a quién sirve tu visibilidad.

\textbf{¿Mi sitio aparece para servir a quien lo busca, o solo para promocionarme yo? ¿A quién dejo afuera con las palabras clave que elijo?}
\end{softbox}
\linead{\linewidth}\\[1mm]
\linead{\linewidth}

\begin{softbox}{milcMostaza}{milcGris}{Estoicismo · Marco Aurelio · cuidado interior}
\textit{``El obstáculo en el camino se vuelve el camino.''}

Los algoritmos de Google cambian todos los meses. Si tu estrategia depende de trucos, se desmorona. Lo que sí depende de ti --- la calidad del contenido, la honestidad de las promesas, la claridad del sitio --- sobrevive los cambios. La paciencia estoica vence al algoritmo.

\textbf{¿Mi SEO depende de trucos que cualquier cambio de algoritmo derriba, o de calidad que aguanta el tiempo? ¿Qué de mi sitio sobrevive a Google?}
\end{softbox}
\linead{\linewidth}\\[1mm]
\linead{\linewidth}

\begin{softbox}{milcTurquesa}{milcGris}{Floridi · lente de la infoesfera}
\textit{``La búsqueda es un acto cooperativo entre el que busca, el que ofrece, y la infoesfera que los conecta.''}

Cada vez que optimizas honestamente, aportas a que la infoesfera funcione mejor para todos. Cada vez que haces SEO de spam (palabras clave irrelevantes, descripciones engañosas), la rompes un poco. Floridi llamaría a esto \emph{degradación informacional}.

\textbf{¿Mi SEO aporta a la cooperación de búsqueda en la infoesfera, o la degrada con ruido? ¿Estoy ayudando a que la gente encuentre lo que necesita?}
\end{softbox}
\linead{\linewidth}\\[1mm]
\linead{\linewidth}

\begin{titlebox}{milcVino}
Compromiso final
\end{titlebox}

\small
\begin{tabularx}{\textwidth}{>{\RaggedRight\arraybackslash}p{3.0cm}YYY}
\rowcolor{milcVino}\color{white}\bfseries Criterio & \color{white}\bfseries Lo logré & \color{white}\bfseries En proceso & \color{white}\bfseries Necesito apoyo\\
Comprensión del tema & Explico ideas centrales con mis palabras. & Reconozco ideas pero falta precisión. & Necesito repasar conceptos base.\\
Pensamiento computacional & Apliqué los 4 pilares al diseñar mi producto. & Apliqué algunos pilares. & Necesito releer la fase 2.\\
Aplicación y producto & Mi producto cumple los criterios y se puede revisar. & Producto funcional con ajustes pendientes. & Aún en construcción.\\
Reflexión 5 dimensiones & Respondí cada dimensión con honestidad. & Respondí algunas con detalle. & Mis respuestas son muy generales.\\
Triángulo de pensamiento & Conecté las 3 voces con mi trabajo real. & Conecté al menos una. & No relaciono las voces con mi proceso.\\
\end{tabularx}
\normalsize

\textbf{Hoy entendí que:}\\[1mm]
\linead{\linewidth}\\[1mm]
\linead{\linewidth}

\textbf{Mi compromiso para la próxima sesión es:}\\[1mm]
\linead{\linewidth}\\[1mm]
\linead{\linewidth}

\begin{softbox}{milcMostaza}{milcGris}{Cita de despedida · Marco Aurelio}
\textit{``El obstáculo en el camino se vuelve el camino. Lo que estorba al camino se vuelve camino.''}\\[1mm]
Cada error de hoy, bien observado, se convierte en pieza del camino que pisarás mañana.
\end{softbox}

\small
\begin{tabularx}{\textwidth}{>{\bfseries}p{3.6cm}Y}
\rowcolor{milcGradeDark!12}Nombre completo: & \linead{10cm}\\
Fecha: & \linead{4cm} \quad Curso/grupo: \linead{4cm}\\
Mi firma: & \linead{9cm}\\
\end{tabularx}
\normalsize

\begin{infoband}{milcGradeDark}
\textbf{Para la próxima semana.} Cada día, abre Google y busca una de tus 3 frases clave estratégicas. Anota en una hoja qué posición ocupa tu sitio y cómo se ve tu resultado (title y description) en el listado. Esa bitácora te enseña a leer cómo te ve Google. El próximo lunes traes la bitácora de 7 días al aula.
\end{infoband}

\end{document}
